%Praeambel, welche den Dokumentenklasse bestimmt, zu ladende Latex-Pakete angibt etc.
%erstmal wählt man den Stil des Dokuments (Schriftgröße, Papierformat etc.)
\documentclass[a4paper,12pt,numbers=noenddot]{scrartcl}
\setlength{\parindent}{1em}
\linespread{1.5} % 1.5-facher Zeilenabstand
% dann lädt man die Pakete, die man benötigt
\usepackage[T1]{fontenc}
\usepackage[ngerman]{babel}
\usepackage{amsmath}
\usepackage{amssymb}
\usepackage{graphicx}
\usepackage{enumitem} % Paket, um einfacher Aufzählungen etc. zu anzupassen
%\usepackage{chapterbib} % added
\usepackage{geometry} % added
\usepackage{lipsum}


\usepackage{textpos,xcolor,rawfonts,latexsym,graphicx,epsf,verbatim} 
\geometry{a4paper, left=35mm, right=25mm, top=20mm, bottom=20mm} % Seiteneinrichtung
\definecolor{fh_orange}{rgb}{0.953,0.201,0}
\definecolor{fh_grau}{rgb}{0.76,0.75,0.76}


% hier beginnt das eigentliche Dokument
\begin{document}
% Globales Setzen der Optionen:

\begin{titlepage}
	%%%%%%%%%%%%%%%%%%%%%%%%%%%%%% -*- Mode: Latex -*- %%%%%%%%%%%%%%%%%%%%%%%%%%%%
	%% 
	%% pa_ba_titelblatt.tex 
	%% 
	%% Copyright (C) 2008 Alexander Sprack / Claudia Holz
	%% 
	%%%%%%%%%%%%%%%%%%%%%%%%%%%%%%%%%%%%%%%%%%%%%%%%%%%%%%%%%%%%%%%%%%%%%%%%%%%%%%%
	\begin{textblock}{6.5}(-1.5,-3)
		\begin{color}{fh_grau}
			\rule{6.8cm}{33cm}    
		\end{color}
	\end{textblock}
	\begin{textblock}{6.5}(-1.2,-0.7)
		% \includegraphics[width=3.8cm]{my-fh-logo}% selbst basteln, falls gewünscht! 
		% Das offizielle Logo ist nicht
		% gestattet!! Bitte BEACHTEN!!!
	\end{textblock}
	\begin{textblock}{6.5}(-1.3,1)
		{\large \textsf{Seminararbeit}}            
	\end{textblock}
	
	\begin{textblock}{7}(4.1,2)
		{\noindent \huge 
			\textsf{\textbf{hier der Titel der Arbeit\\[0.3cm] 
					\Large  Zweizeiliger Untertitel\\[0.05cm]
					sofern vorhanden}} }
	\end{textblock}
	
	
	\begin{textblock}{6}(4.1,6.5)\noindent
		\textsf{An der Fachhochschule Dortmund\\
			im Fachbereich Informatik\\
			Studiengang Informatik \\
			erstellte Seminararbeit \\
			für das BIP:\\
			GenAI meets Reality}
	\end{textblock}
	
	\begin{textblock}{6.5}(-0.9,10.8)
		\noindent
		\textsf{von \\
			Vorname Name \\
			Matr.-Nr. xxxx xxx\\[0.7cm]
			Betreuerin: Prof. Dr. Klaus Kaiser \\[0.5cm]
			Dortmund, \today}    
	\end{textblock}
	
\end{titlepage}





\begin{abstract}
    \lipsum[1-2]
\end{abstract}

\newpage
\tableofcontents

\part{Projektbericht}

\section{Einleitung}
    Eine Einleitung in die Problemstellung und Lösungskonzepten. Verwenden Sie bitte Zitate und Referenzieren Sie alle eingefügten Grafiken im Text.

    RAG zum Beispiel wurde eingefügt in der Arbeit von Lewis et al. \cite{rag} und eine Darstellung des Antwortprozesses ist gegeben in Abbildung \ref{fig:rag}.

    \begin{figure}
        \begin{center}
            \includegraphics[scale=0.6]{images/rag.png}
        \end{center}
        \caption{Darstellung des Antwortprozesses in einem RAG System.}\label{fig:rag}
    \end{figure}

\section{Daten}
    \lipsum[1-2]

\section{Umsetzung}
    \lipsum[1-2]

\section{Ergebnis}
    \lipsum[1-2]

\section{Fazit und Ausblick}
    \lipsum[1-2]

\begin{thebibliography}{9}
\bibitem{rag}
Lewis, Patrick, et al. ``Retrieval-augmented generation for knowledge-intensive nlp tasks.'' \textit{Advances in neural information processing systems 33} (2020): 9459-9474.
\end{thebibliography}


\part{Selbstreflektion}
    \lipsum[1-2]

\end{document}          

